% ====================================================================
% FF3-7  준안정 vs 불안정의 두 동역학 경로 (Appendix 문맥, §A.6 제안)
%   좌표 = (a) eq:app-cnt·eq:app-rstar 의 무차원화 ΔG/ΔG*=3(r/r*)²−2(r/r*)³,
%          (b) eq:app-ch-R 의 무차원화 R/R(k_m)=4(k/k_c)²[1−(k/k_c)²] —
%   식 그대로의 수치 평가 (gen_coords_ff3.py).
%   사용 라이브러리: arrows.meta, calc 만 (appendix preamble 내).
% ====================================================================
\begin{figure}[t]
\centering
\begin{tikzpicture}
% =============== (a) classical nucleation barrier ===============
\begin{scope}[x=3.5cm,y=1.5cm]
\draw[-{Latex[length=1.5mm]}] (0,-1.30) -- (0,1.75)
  node[above,font=\scriptsize] {$\Delta G/\Delta G^\ast$};
\draw[-{Latex[length=1.5mm]}] (0,0) -- (1.75,0)
  node[below,font=\scriptsize] {$r/r^\ast$};
\foreach \xx in {0.5,1.0,1.5} {\draw (\xx,0.02) -- (\xx,-0.02) node[below,font=\scriptsize] {\xx};}
\foreach \yy in {-1,1} {\draw (-0.012,\yy) -- (0.012,\yy) node[left,font=\scriptsize] {$\yy$};}
% surface cost +3x^2 (thin dashed, clipped)
\draw[densely dashed,black!55] plot[smooth] coordinates {(0,0) (0.1,0.03) (0.2,0.12) (0.3,0.27) (0.4,0.48) (0.5,0.75) (0.6,1.08) (0.7,1.47) (0.72,1.555)};
\node[black!55,font=\tiny,anchor=west] at (0.76,1.50) {표면 비용 $+3(r/r^\ast)^2$};
% volume gain -2x^3 (thin dotted, clipped)
\draw[densely dotted,black!55] plot[smooth] coordinates {(0,0) (0.15,-0.00675) (0.3,-0.054) (0.45,-0.1823) (0.6,-0.432) (0.7,-0.686) (0.8,-1.024) (0.85,-1.228)};
\node[black!55,font=\tiny,anchor=west] at (0.88,-1.05) {부피 이득 $-2(r/r^\ast)^3$};
% total (thick)
\draw[very thick] plot[smooth] coordinates {(0,0) (0.05,0.00725) (0.1,0.028) (0.15,0.06075) (0.2,0.104) (0.25,0.1563) (0.3,0.216) (0.35,0.2818) (0.4,0.352) (0.45,0.4253) (0.5,0.5) (0.55,0.5748) (0.6,0.648) (0.65,0.7183) (0.7,0.784) (0.75,0.8438) (0.8,0.896) (0.85,0.9393) (0.9,0.972) (0.95,0.9928) (1,1) (1.05,0.9928) (1.1,0.968) (1.15,0.9223) (1.2,0.852) (1.25,0.7538) (1.3,0.624) (1.35,0.4593) (1.4,0.256) (1.45,0.01087) (1.5,-0.28) (1.55,-0.6198) (1.6,-1.012) (1.63,-1.269)};
% critical point
\fill (1,1) circle (1.2pt);
\draw[densely dotted] (1,0) -- (1,1);
\draw[densely dotted] (0,1) -- (1,1);
\node[font=\tiny,anchor=south west] at (1.01,1.02) {$(r^\ast,\ \Delta G^\ast)$};
\node[font=\tiny,anchor=west,align=left] at (0.06,-0.55)
  {$r<r^\ast$: 계면 비용 지배\\$\to$ 핵이 녹아 되돌아감};
\node[font=\tiny,anchor=west,align=left] at (1.12,-0.30)
  {$r>r^\ast$:\\성장};
\node[font=\scriptsize,anchor=west] at (0.06,1.70) {(a) 준안정 --- 핵생성 장벽};
\end{scope}
% =============== (b) spinodal decomposition growth band ===============
\begin{scope}[xshift=7.6cm,x=4.6cm,y=2.35cm]
\draw[-{Latex[length=1.5mm]}] (0,-0.62) -- (0,1.22)
  node[above,font=\scriptsize] {$R(k)/R(k_m)$};
\draw[-{Latex[length=1.5mm]}] (0,0) -- (1.22,0)
  node[below,font=\scriptsize] {$k/k_c$};
\foreach \xx in {0.5,1.0} {\draw (\xx,0.013) -- (\xx,-0.013) node[below,font=\scriptsize] {\xx};}
\foreach \yy in {1} {\draw (-0.008,\yy) -- (0.008,\yy) node[left,font=\scriptsize] {$\yy$};}
% growth rate curve
\draw[very thick] plot[smooth] coordinates {(0,0) (0.1,0.0396) (0.2,0.1536) (0.3,0.3276) (0.4,0.5376) (0.5,0.75) (0.6,0.9216) (0.7071,1.0) (0.8,0.9216) (0.9,0.6156) (1.0,0) (1.05,-0.452)};
% k_m marker
\fill (0.7071,1.0) circle (1.2pt);
\draw[densely dotted] (0.7071,0) -- (0.7071,1.0);
\node[font=\tiny,anchor=west] at (0.735,1.04) {$k_m=k_c/\sqrt2$};
\node[font=\tiny,anchor=west,align=left] at (0.06,0.86) {성장 밴드\\$0<k<k_c$ ($R>0$)};
\node[font=\tiny,anchor=west,align=left] at (1.07,-0.28) {$k>k_c$: 감쇠\\(구배 벌칙)};
\node[font=\scriptsize,anchor=west] at (0.03,1.15) {(b) 불안정 --- spinodal 분해};
\end{scope}
\end{tikzpicture}
\caption{준안정과 불안정 영역의 두 동역학 경로 --- 좌표는 식 그대로의 수치 평가다.
(a) 고전 핵생성: 구형 핵의 자유에너지 식~\eqref{eq:app-cnt} 을 임계값
식~\eqref{eq:app-rstar} 로 무차원화하면 $\Delta G/\Delta G^\ast=3(r/r^\ast)^2-2(r/r^\ast)^3$ 이다.
표면 비용($+r^2$, 파선)과 부피 이득($-r^3$, 점선)의 경쟁이 $r^\ast$ 에서 장벽 $\Delta G^\ast$ 를
만들고, $r<r^\ast$ 인 핵은 되돌아가며 $r>r^\ast$ 만 성장한다 --- 준안정 영역의 분해가
$\exp(-\Delta G^\ast/k_BT)$ 로 지수 억제되는 까닭이다. (b) spinodal 분해: 불안정 영역($f''<0$)의
모드별 성장률 식~\eqref{eq:app-ch-R} 을 최대 성장 모드로 무차원화하면
$R(k)/R(k_m)=4(k/k_c)^2[1-(k/k_c)^2]$ 이며, 장벽 없이 $0<k<k_c$ 밴드 전체가 지수 성장하고
$k_m=k_c/\sqrt2$ 물결이 가장 빨리 자란다. 장벽 유무가 두 영역(binodal--spinodal 사이 vs spinodal
안쪽)의 조작적 차이다.}
\label{fig:cand-ff3-7}
\end{figure}
